\documentclass[10pt]{article}
\usepackage[margin=0.8in]{geometry}
\usepackage{amsmath,amssymb,hyperref}
\hypersetup{colorlinks=true,urlcolor=blue,linkcolor=blue}
\setlength{\parindent}{0pt}\setlength{\parskip}{4pt}
\begin{document}
\begin{center}
{\Large\bfseries Problem 00000000358: the normalized error is bounded}

lizaixi01 \quad AI-led research and formalization \quad 23 October 2026
\end{center}
\textbf{Result and scope.} The printed statement assigns amplitude at least $t^{11}$ to the error of $N(t)/\operatorname{vol}(tB)$ for the Leech lattice. We prove that this normalized quantity is bounded for $t\ge1$, hence its deviation from any fixed real center cannot have cofinal amplitude $c\,t^{11}$ for any $c>0$. Only this necessary amplitude clause is refuted. No assertion is made that sign changes never occur, or that the same claim for an \emph{unnormalized} discrepancy is false.

The source leaves $B$ unspecified. Our theorem covers every fixed bounded measurable body with positive finite Lebesgue volume, including the ordinary Euclidean unit ball and convex bodies with interior. It does not silently substitute a growing or zero-volume body. The proof treats the full standard coordinate recipe with a binary code as input, universally in that input; the actual extended Golay code is therefore a particular instance. It does not call arbitrary codes Golay or reprove the classification or existence theory of the Leech lattice.

\section*{1. The actual coordinate construction}
Wilson's author lecture notes \cite{W}, printed p.~2, Section~3, define the integer-coordinate model using the following three conditions on $z\in\mathbb Z^{24}$, for some $m\in\mathbb Z$:
\[
 z_i\equiv m\pmod2,\qquad
 \left(\frac{z_i-m}{2}\bmod2\right)_{i=1}^{24}\in C,\qquad
 \sum_i z_i\equiv4m\pmod8.
\]
Here $C$ is the extended binary Golay code. Printed p.~3 describes the norm divided by $8$, corresponding to multiplying ordinary Euclidean coordinates by $1/\sqrt8$.

Lean implements all three conditions using actual integers, remainders, and $\mathbb Z/2\mathbb Z$, then forms their actual image in $E=\mathbb R^{24}$ under
\[
 P_a(z)=(az_i)_{i=1}^{24},\qquad a>0.
\]
It proves this coordinate map injective. The construction is denoted $\Lambda_{a,C}$. The formal theorem holds for every set $C\subseteq(\mathbb Z/2\mathbb Z)^{24}$ and every fixed $a>0$, so it applies to the genuine Golay choice with $a=1/\sqrt8$ and also to fixed alternative scalings. The only property used in the estimate is proved from this full construction: every point has a preimage in the scaled integer grid. This is not an unproved lattice-containment hypothesis.

Define the actual point set and count
\[
 S_t=\Lambda_{a,C}\cap(t\cdot B),\qquad N(t)=\#S_t.
\]
All geometry uses Mathlib's actual Euclidean space and its Lebesgue volume. We first prove $S_t$ finite for every real $t$; the natural cardinality never relies on the default value for an infinite set.

\section*{2. A uniform coordinate count}
Boundedness gives a real $R>0$ with $\|x\|\le R$ for every $x\in B$. If $P_a(z)\in S_t$, the actual norm and scalar-image identities imply
\[
 |z_i|\le |t|R/a\quad\hbox{for every }i.
\]
Let $Q=\lceil |t|R/a\rceil\in\mathbb N$. Thus $S_t$ is contained in the image of the finite integer box $\{-Q,\ldots,Q\}^{24}$. Lean proves both this inclusion and the box's exact cardinality $(2Q+1)^{24}$, yielding
\[
 N(t)\le(2\lceil |t|R/a\rceil+1)^{24}.
\]
For $t\ge1$, the actual ceiling bound gives
\[
 N(t)\le(2R/a+3)^{24}t^{24}.
\]
No count is defined to equal this bound: it is an inequality for the genuine finite intersection.

\section*{3. Actual volume and the cofinal contradiction}
Put $v=\operatorname{vol}(B)\in(0,\infty)$. The actual Haar/Lebesgue scalar-image theorem in dimension $24$ gives
\[
 \operatorname{vol}(t\cdot B)=t^{24}v\qquad(t\ge0).
\]
The Lean proof uses the real value of the actual extended-nonnegative measure, explicitly proving it positive and finite where division occurs. Consequently, with $M=(2R/a+3)^{24}/v>0$,
\[
 0\le\frac{N(t)}{\operatorname{vol}(tB)}\le M\qquad(t\ge1).
\]
For any fixed real center $b$, the absolute deviation is at most $M+|b|$. This includes whatever constant main term is used for the normalized lattice density; no determinant or covolume value is assumed.

The formal amplitude condition is the weak, cofinal version
\[
 \forall T\in\mathbb R\ \exists t\ge T:\quad
 c\,t^{11}\le\left|\frac{N(t)}{\operatorname{vol}(tB)}-b\right|.
\]
For any $c>0$, take $T=\max\{1,(M+|b|)/c+1\}$. Then $t\ge T$ implies $c\,t^{11}\ge ct>M+|b|$, contradicting the actual upper bound. This disproves even the cofinal condition, so also any stronger eventual amplitude lower bound. The theorem quantifies over every fixed center and positive multiplier, not a single sample radius.

An explicit final instance uses the genuine Euclidean closed unit ball, proving its measurability, boundedness, positive finite volume, and the positive scaling $1/\sqrt8$. The universal code theorem includes the extended Golay input without assigning it an unproved concrete generator matrix.

\section*{4. Verification and attribution}
All 28 named declarations (16 theorems plus definitions/abbreviations) are axiom-audited. The package preserves the complete bilingual source, this LaTeX/PDF, a pinned Lean project, source-bound build and fresh-kernel logs, and SHA256 manifests. No \texttt{sorry}, \texttt{native\_decide} or additional axiom is used. Worker and separate supervising AI inspect the proof scope and all rendered pages. No organizer acceptance or human verification is asserted.

The elementary grid estimate is proved directly, not assumed from an asymptotic lattice-counting theorem. Wilson is cited only for the actual coordinate recipe and scaling. This selectively produced result is not appended to the frozen CPU or workflow experiments.
\begin{thebibliography}{1}
\bibitem{W} R.~A.~Wilson, \emph{The Leech lattice}, QMUL Pure Mathematics Seminar, 8 December 2008, \href{https://webspace.maths.qmul.ac.uk/r.a.wilson/talks_files/Leech.pdf}{author notes}, printed pp.~2--3, Section~3 and the norm normalization in Section~4.
\end{thebibliography}
\end{document}
